% Lecture 5 notes: \input by notes.tex, \subimport'ed by the semester book.
% Exam questions (экзамен): 10 (основная теорема арифметики, задача факторизации, перебор
% делителей), 11 (алгоритм Евклида, теорема Эйлера, RSA), 15 (алгоритм p - 1 Полларда).
\begin{lecture}{5}{Задача факторизации, RSA и алгоритм \texorpdfstring{$p - 1$}{p - 1} Полларда}

\section{Криптосистема с открытым ключом, основанная на факторизации}

В предыдущих лекциях стойкость протоколов (Диффи"--~Хеллман, Эль-Гамаль) опиралась на
сложность задачи дискретного логарифмирования. В 1977 г. Р.~Ривест, А.~Шамир и Л.~Адлеман
предложили криптосистему RSA~\cite{rsa}, стойкость которой связана с другой классической
задачей "--- \emph{разложением целого числа на простые множители} (факторизацией). Её
устройство очень простое:
\begin{itemize}
  \item Алиса выбирает два больших простых $p$ и $q$ (сейчас "--- по $1024$ бита и более) и
    публикует их произведение $N = pq$ вместе с показателем $e$;
  \item шифрование и расшифрование "--- возведение в степень по модулю $N$, т.\,е.
    $\bigO(\log N)$ умножений (лекция~6 первого семестра);
  \item чтобы расшифровывать, нужно знать $\varphi(N) = (p - 1)(q - 1)$, а для этого "---
    разложение $N$.
\end{itemize}
Перемножить два простых числа легко, а восстановить их по произведению "--- по-видимому,
трудно: лучшие известные алгоритмы субэкспоненциальны, и рекорд факторизации модуля вида RSA
(RSA-250, $829$ бит, 2020 г.) потребовал тысяч процессоро-лет. В этой лекции мы докажем
основную теорему арифметики, сформулируем задачу факторизации, построим RSA и выясним, какие
сведения о ключе равносильны разложению $N$. В конце разберём первый «специальный» алгоритм
факторизации "--- метод $p - 1$ Полларда, который объясняет, почему простые для RSA
выбирают с осторожностью.

\section{Основная теорема арифметики}

Натуральное число $p > 1$ называется \emph{простым}, если у него нет делителей, кроме $1$ и
$p$, и \emph{составным} в противном случае. Число $1$ не является ни простым, ни составным.

\begin{lemma}\label{l5:least-divisor}
Наименьший делитель $d > 1$ натурального числа $n > 1$ "--- простое число. Если $n$ составное,
то $d \le \sqrt n$.
\end{lemma}

\begin{proof}
Если бы $d$ имело делитель $1 < d' < d$, то $d' \mid n$ вопреки минимальности $d$. Если $n$
составное, то $n = d \cdot n'$, где $n' > 1$ "--- тоже делитель $n$, поэтому $n' \ge d$ и
$n = d n' \ge d^2$.
\end{proof}

\begin{lemma}[Евклид]\label{l5:euclid-lemma}
Если $p$ простое и $p \mid ab$, то $p \mid a$ или $p \mid b$.
\end{lemma}

\begin{proof}
Пусть $p \nmid a$. Так как делители $p$ "--- только $1$ и $p$, то $\gcd(a, p) = 1$, и по
соотношению Безу (теорема~\ref{l5:bezout} ниже) $1 = ua + vp$ для некоторых целых $u, v$.
Умножая на $b$, получаем $b = u \cdot ab + vb \cdot p$; оба слагаемых делятся на $p$,
значит, $p \mid b$.
\end{proof}

Индукцией по числу сомножителей: если $p \mid a_1 a_2 \cdots a_k$, то $p$ делит одно из
$a_i$; в частности, если $p \mid q_1 \cdots q_k$, где $q_i$ простые, то $p = q_i$ для
некоторого $i$.

\begin{theorem}[основная теорема арифметики]\label{l5:fta}
Каждое натуральное $n > 1$ раскладывается в произведение простых чисел
\[
  n = p_1 p_2 \cdots p_k,
\]
и это разложение единственно с точностью до порядка сомножителей. Группируя одинаковые
множители, получаем \emph{каноническое разложение}
$n = q_1^{\alpha_1} q_2^{\alpha_2} \cdots q_s^{\alpha_s}$, $q_1 < q_2 < \cdots < q_s$,
$\alpha_i \ge 1$.
\end{theorem}

\begin{proof}
\emph{Существование} "--- индукция по $n$. Если $n$ простое, разложение состоит из одного
множителя. Иначе по лемме~\ref{l5:least-divisor} $n = p \cdot n'$, где $p$ простое и
$1 < n' < n$; по предположению индукции $n'$ "--- произведение простых, значит, и $n$ тоже.

\emph{Единственность} "--- индукция по $n$. Пусть $n = p_1 \cdots p_k = q_1 \cdots q_l$ "---
два разложения. Так как $p_1 \mid q_1 \cdots q_l$, по лемме~\ref{l5:euclid-lemma}
$p_1 = q_j$ для некоторого $j$. Сокращая на $p_1$, получаем два разложения числа
$n / p_1 < n$; если $n/p_1 = 1$, то $k = l = 1$, иначе по предположению индукции они
совпадают с точностью до порядка, а значит, совпадают и исходные разложения.
\end{proof}

\begin{remark}
Единственность "--- нетривиальное свойство: в кольце $\ZZ[\sqrt{-5}]$
\[
  6 = 2 \cdot 3 = (1 + \sqrt{-5})(1 - \sqrt{-5}),
\]
и все четыре множителя неразложимы. Ключевым в доказательстве было соотношение Безу, т.\,е.
алгоритм Евклида.
\end{remark}

\checkyourself{Где в доказательстве единственности используется простота $p_1$? Приведите
пример, показывающий, что лемма Евклида неверна для составного $p$.}

\section{Задача факторизации}

\begin{definition}
\emph{Задача факторизации}: по данному натуральному $N > 1$ найти его каноническое разложение.
\end{definition}

Достаточно научиться находить \emph{какой-нибудь} нетривиальный делитель $1 < d < N$
составного числа: тогда задача сводится к разложению $d$ и $N/d$, а число таких шагов не
больше $\log_2 N$ (столько простых множителей может быть у $N$). Проверить, что число
простое, можно быстро "--- за полиномиальное время вероятностными тестами Миллера"--~Рабина и
Соловея"--~Штрассена (лекция~8 первого семестра). Поэтому всюду ниже $N$ "--- заведомо
составное число, а в криптографических приложениях обычно $N = pq$ "--- произведение двух
различных простых сравнимой длины.

\subsection{Перебор делителей}

Лемма~\ref{l5:least-divisor} сразу даёт алгоритм.

\begin{algorithm}[перебор делителей]\label{l5:trial}
\begin{algorithmic}[1]
\Require составное $N$.
\Ensure нетривиальный делитель $N$.
\For{$d = 2, 3, \dots, \floor{\sqrt N}$}
  \If{$d \mid N$} \Return $d$ \EndIf
\EndFor
\end{algorithmic}
\end{algorithm}

\begin{proposition}\label{l5:trial-complexity}
Перебор делителей находит наименьший простой делитель $p$ числа $N$ за $p - 1 \le \sqrt N$
делений. Для $N = pq$ с простыми $p \approx q$ это $\Theta(\sqrt N) = \Theta(2^{n/2})$ делений,
где $n = \log_2 N$ "--- длина входа: сложность \emph{экспоненциальна} в длине входа.
\end{proposition}

\begin{proof}
Первый найденный делитель "--- наименьший делитель $p > 1$; он прост и не превосходит
$\sqrt N$ по лемме~\ref{l5:least-divisor}. Каждое деление $n$-битного числа на число длины
не более $n/2$ стоит $\bigO(n^2)$ битовых операций, так что общая сложность
$\bigO(2^{n/2} n^2)$, а при $p \approx q \approx \sqrt N$ перебор не может закончиться раньше.
\end{proof}

Если делить только на простые числа, число делений уменьшается до
$\pi(\sqrt N) \sim 2\sqrt N / \ln N$, что по-прежнему экспоненциально. Для $2048$-битного
модуля RSA это около $10^{305}$ делений "--- задача, безнадёжная при любых мыслимых
ресурсах. Тем не менее перебор полезен как первый шаг любого метода: он быстро отделяет
\emph{малые} простые множители.

Сложность алгоритмов факторизации удобно сравнивать с помощью функции
\[
  L_N[\alpha, c] = \exp\bigl(c\,(\ln N)^{\alpha} (\ln\ln N)^{1 - \alpha}\bigr), \qquad
  0 \le \alpha \le 1.
\]
При $\alpha = 0$ это $(\ln N)^c$ "--- полиномиальная сложность, при $\alpha = 1$ "--- $N^c$,
экспоненциальная (перебор делителей "--- $L_N[1, \frac12]$); промежуточные значения
называются \emph{субэкспоненциальными}. Квадратичное решето (лекция~6) работает за
$L_N[\frac12, 1]$, метод эллиптических кривых Ленстры (лекция~10) "--- за
$L_p[\frac12, \sqrt 2]$, где $p$ "--- искомый делитель, а лучший известный алгоритм, решето
числового поля, "--- за $L_N[\frac13, \sqrt[3]{64/9}]$, $\sqrt[3]{64/9} \approx 1{,}923$.
Полиномиального алгоритма факторизации на классическом компьютере не известно (на квантовом
его даёт алгоритм Шора).

\checkyourself{Почему для доказательства того, что $N$ составное, не нужно знать ни одного
его делителя? Сколько делений выполнит перебор для $N = 8051$, и почему этот пример плох для
перебора, но хорош для метода Ферма из лекции~6 ($8051 = 90^2 - 7^2$)?}

\section{Алгоритм Евклида и теорема Эйлера}

\subsection{Основной и расширенный алгоритм Евклида}

Напомним результаты лекции~3 первого семестра. Для целых $a \ge b \ge 0$, не равных нулю
одновременно, $\gcd(a, b) = \gcd(b, a \bmod b)$ (общие делители пар $(a, b)$ и
$(b, a - qb)$ совпадают), и $\gcd(a, 0) = a$.

\begin{algorithm}[расширенный алгоритм Евклида]\label{l5:xgcd}
\begin{algorithmic}[1]
\Require целые $a \ge b \ge 0$, $a > 0$.
\Ensure $g = \gcd(a, b)$ и $u, v$: $ua + vb = g$.
\State $(r_0, u_0, v_0) \gets (a, 1, 0)$, \quad $(r_1, u_1, v_1) \gets (b, 0, 1)$, \quad $i \gets 1$
\While{$r_i \ne 0$}
  \State $q_i \gets \floor{r_{i-1} / r_i}$
  \State $(r_{i+1}, u_{i+1}, v_{i+1}) \gets (r_{i-1}, u_{i-1}, v_{i-1}) - q_i (r_i, u_i, v_i)$,
    \quad $i \gets i + 1$
\EndWhile
\State \Return $(r_{i-1}, u_{i-1}, v_{i-1})$
\end{algorithmic}
\end{algorithm}

Без столбцов $u, v$ это \emph{основной} алгоритм Евклида: $r_{i+1} = r_{i-1} \bmod r_i$.

\begin{theorem}[соотношение Безу]\label{l5:bezout}
Расширенный алгоритм Евклида корректен: на каждом шаге $r_i = u_i a + v_i b$, последний ненулевой
остаток равен $\gcd(a, b)$. В частности, $\gcd(a, b) = ua + vb$ для некоторых целых $u, v$.
Число делений не превосходит $2\log_2 b + 1$.
\end{theorem}

\begin{proof}
Равенство $r_i = u_i a + v_i b$ верно для $i = 0, 1$ и сохраняется, так как тройки
$(r, u, v)$ преобразуются одной и той же линейной операцией. Остатки строго убывают, поэтому
алгоритм останавливается, а $\gcd(r_{i-1}, r_i)$ не меняется по шагам и в конце равен
$\gcd(r_{i-1}, 0) = r_{i-1}$.

Оценка: покажем, что $r_{i+1} < r_{i-1} / 2$. Если $r_i \le r_{i-1}/2$, то
$r_{i+1} < r_i \le r_{i-1}/2$; иначе $q_i = 1$ и $r_{i+1} = r_{i-1} - r_i < r_{i-1}/2$. Значит,
за каждые два шага остаток уменьшается более чем вдвое, начиная с $r_1 = b$.
\end{proof}

Итак, алгоритм Евклида выполняет $\bigO(n)$ делений $n$-битных чисел, т.\,е. $\bigO(n^3)$
битовых операций (аккуратный анализ даёт $\bigO(n^2)$). Если $\gcd(a, m) = 1$, то из
$ua + vm = 1$ получаем $a^{-1} \equiv u \pmod m$: обратный элемент вычисляется так же быстро.

\begin{example}\label{l5:xgcd-example}
$17^{-1} \bmod 3120$: расширенный алгоритм Евклида для $(3120, 17)$ даёт
\[
  \xgcdtable{3120}{17}
  \qquad 2 \cdot 3120 - 367 \cdot 17 = 1,
\]
поэтому $17^{-1} \equiv -367 \equiv 2753 \pmod{3120}$.
\end{example}

\subsection{Функция и теорема Эйлера}

\begin{definition}
\emph{Функция Эйлера} $\varphi(N)$ "--- число целых $0 \le a < N$, взаимно простых с $N$,
т.\,е. порядок группы обратимых элементов $\Znx{N}$.
\end{definition}

\begin{proposition}\label{l5:phi}
Если $p$ простое, то $\varphi(p) = p - 1$ и $\varphi(p^k) = p^k - p^{k-1}$. Если $p \ne q$
простые, то $\varphi(pq) = (p - 1)(q - 1)$. В общем случае $\varphi$ мультипликативна:
$\varphi(mn) = \varphi(m)\varphi(n)$ при $\gcd(m, n) = 1$, и
$\varphi(N) = N \prod_{p \mid N} (1 - 1/p)$.
\end{proposition}

\begin{proof}
Не взаимно просты с $p^k$ ровно числа, кратные $p$: их $p^{k-1}$ среди $0, \dots, p^k - 1$.
Не взаимно просты с $pq$ числа, кратные $p$ ($q$ штук) или $q$ ($p$ штук); число $0$ кратно
обоим, а других общих кратных на отрезке $[0, pq)$ нет. Поэтому
$\varphi(pq) = pq - p - q + 1 = (p - 1)(q - 1)$. Мультипликативность следует из китайской
теоремы об остатках (лекция~2): изоморфизм $\Zn{mn} \cong \Zn{m} \times \Zn{n}$ переводит
обратимые элементы в пары обратимых, т.\,е. $\Znx{mn} \cong \Znx{m} \times \Znx{n}$.
Последняя формула получается перемножением формул для степеней простых.
\end{proof}

\begin{theorem}[Эйлер]\label{l5:euler}
Если $\gcd(a, N) = 1$, то $a^{\varphi(N)} \equiv 1 \pmod N$.
\end{theorem}

\begin{proof}
Пусть $x_1, \dots, x_k$, $k = \varphi(N)$, "--- все вычеты, взаимно простые с $N$. Числа
$a x_1, \dots, a x_k$ тоже взаимно просты с $N$ и попарно несравнимы: из
$a x_i \equiv a x_j$ после умножения на $a^{-1}$ следует $x_i \equiv x_j$. Значит, это те же
вычеты, переставленные, и
\[
  \prod_{i=1}^{k} (a x_i) \equiv \prod_{i=1}^{k} x_i \pmod N, \qquad\text{т.\,е.}\qquad
  a^{k} \prod x_i \equiv \prod x_i \pmod N.
\]
Произведение $\prod x_i$ обратимо по модулю $N$; сокращая на него, получаем $a^k \equiv 1$.
(Иначе: порядок элемента $a$ делит порядок группы $\Znx{N}$ по теореме Лагранжа.)
\end{proof}

\begin{corollary}[малая теорема Ферма]\label{l5:fermat}
Если $p$ простое и $p \nmid a$, то $a^{p-1} \equiv 1 \pmod p$; для любого $a$
$a^p \equiv a \pmod p$.
\end{corollary}

\begin{corollary}\label{l5:exp-reduce}
Если $\gcd(a, N) = 1$ и $k \equiv l \pmod{\varphi(N)}$, то $a^k \equiv a^l \pmod N$: показатели
можно приводить по модулю $\varphi(N)$.
\end{corollary}

\checkyourself{Вычислите $\varphi(100)$ и $7^{2026} \bmod 100$. Верна ли теорема Эйлера при
$\gcd(a, N) > 1$?}

\section{Криптосистема RSA}

\begin{algorithm}[RSA]\label{l5:rsa}
\begin{enumerate}
  \item \textbf{[Генерация ключей]} Алиса выбирает различные простые $p$, $q$ (случайные,
    одинаковой длины), вычисляет $N = pq$, $\varphi(N) = (p - 1)(q - 1)$, выбирает открытую
    экспоненту $e$, $1 < e < \varphi(N)$, $\gcd(e, \varphi(N)) = 1$ (часто $e = 65537$), и
    вычисляет $d \equiv e^{-1} \pmod{\varphi(N)}$ расширенным алгоритмом Евклида.
    \emph{Открытый ключ} "--- $(N, e)$, \emph{закрытый} "--- $d$ (а также $p$, $q$).
  \item \textbf{[Шифрование]} Боб представляет сообщение числом $0 \le m < N$ и отправляет
    $c = m^e \bmod N$.
  \item \textbf{[Расшифрование]} Алиса вычисляет $m = c^d \bmod N$.
\end{enumerate}
\end{algorithm}

Условие $\gcd(e, \varphi(N)) = 1$ необходимо: иначе $d$ не существует, а отображение
$x \mapsto x^e$ не инъективно на $\Znx{N}$. Например, если простое $r$ делит
$\gcd(e, p - 1)$, то в циклической группе $\Znx{p}$ есть элемент $\zeta \ne 1$ порядка $r$, и
$x$, $\zeta x$ имеют одинаковые $e$-е степени по модулю $p$.

\begin{theorem}[корректность RSA]\label{l5:rsa-correct}
Если $ed \equiv 1 \pmod{\varphi(N)}$, $N = pq$, $p \ne q$ простые, то $m^{ed} \equiv m \pmod N$
для \emph{всех} целых $m$, в том числе не взаимно простых с $N$.
\end{theorem}

\begin{proof}
Пусть $ed = 1 + k(p - 1)(q - 1)$, $k \ge 0$. Докажем, что $m^{ed} \equiv m \pmod p$. Если
$p \mid m$, обе части сравнимы с нулём. Иначе по малой теореме Ферма
\[
  m^{ed} = m \cdot \bigl(m^{p-1}\bigr)^{k(q-1)} \equiv m \cdot 1 = m \pmod p.
\]
Точно так же $m^{ed} \equiv m \pmod q$. Итак, $p$ и $q$ делят $m^{ed} - m$, а так как они
различны и просты, то и $N = pq$ делит $m^{ed} - m$ (это частный случай китайской теоремы об
остатках: система $x \equiv m \pmod p$, $x \equiv m \pmod q$ имеет единственное решение по
модулю $N$).
\end{proof}

Если бы мы опирались только на теорему Эйлера, случай $\gcd(m, N) > 1$ остался бы
непокрытым; именно поэтому доказательство ведётся отдельно по модулям $p$ и $q$.
На практике такие $m$ встречаются с вероятностью $1 - \varphi(N)/N = 1/p + 1/q - 1/N \approx 2/\sqrt N$ (при $p \approx q$), и
наткнуться на одно из них означает разложить $N$ (через $\gcd(m, N)$).

\begin{remark}
Из доказательства видно, что достаточно $ed \equiv 1 \pmod{\lambda(N)}$, где
$\lambda(N) = \lcm(p - 1, q - 1)$ "--- \emph{функция Кармайкла} (показатель группы $\Znx{N}$):
стандарт PKCS\,\#1 допускает $d = e^{-1} \bmod \lambda(N)$, это число меньше.
\end{remark}

\begin{example}\label{l5:rsa-hand}
$p = 61$, $q = 53$: $N = 3233$, $\varphi(N) = 60 \cdot 52 = 3120 = 2^4 \cdot 3 \cdot 5 \cdot 13$.
Возьмём $e = 17$; по примеру~\ref{l5:xgcd-example} $d = 2753$. Шифрование $m = 65$, $17 = 10001_2$:
\[
  65 \xrightarrow{\,x^2\,} 992 \xrightarrow{\,x^2\,} 1232 \xrightarrow{\,x^2\,} 1547
  \xrightarrow{\,x^2\,} 789 \xrightarrow{\,\cdot 65\,} 2790 = c.
\]
Расшифрование $2790^{2753} \bmod 3233$ возвращает $65$ (проверьте!).
\end{example}

\subsection{Расшифрование с помощью китайской теоремы об остатках}

Закрытый показатель $d$ имеет ту же длину, что и $N$, поэтому расшифрование дороже
шифрования с малым $e$. Владелец ключа знает $p$ и $q$ и может работать по модулям $p$ и $q$
отдельно:
\[
  d_p = d \bmod (p - 1), \quad d_q = d \bmod (q - 1), \quad
  m_p = c^{d_p} \bmod p, \quad m_q = c^{d_q} \bmod q,
\]
после чего $m$ восстанавливается по КТО, например по формуле Гарнера
\[
  m = m_q + q \cdot \bigl((m_p - m_q)\, q^{-1} \bmod p\bigr).
\]
Корректность: $m \equiv c^d \pmod p$, а $c^d \equiv c^{d_p} \pmod p$ при $p \nmid c$ по малой
теореме Ферма ($d - d_p$ кратно $p - 1$), при $p \mid c$ обе части равны нулю. Модули и
показатели вдвое короче; так как возведение в степень стоит $\bigO(n^3)$ при школьном умножении,
каждое из двух возведений в $8$ раз дешевле, а в сумме выигрыш около $4$ раз.

\begin{example}
Для примера~\ref{l5:rsa-hand}: $d_p = 2753 \bmod 60 = 53$, $d_q = 2753 \bmod 52 = 49$,
$c_p = 2790 \bmod 61 = 45$, $c_q = 2790 \bmod 53 = 34$; $m_p = 45^{53} \bmod 61 = 4$,
$m_q = 34^{49} \bmod 53 = 12$. Далее $53^{-1} \bmod 61 = 38$,
$(4 - 12) \cdot 38 \bmod 61 = 1$, и $m = 12 + 53 \cdot 1 = 65$.
\end{example}

\taskexample{rsa-toy}

\begin{proposition}[сложность RSA]
Пусть $n = \log_2 N$. Шифрование с фиксированным малым $e$ стоит $\bigO(n^2)$, расшифрование "---
$\bigO(n^3)$ битовых операций (школьное умножение). Генерация ключей требует в среднем
$\bigO(n)$ проверок простоты случайных $n/2$-битных чисел (по асимптотическому закону
распределения простых доля простых среди них $\approx 2/(n \ln 2)$), т.\,е. $\bigO(n^4)$
операций при тесте Миллера"--~Рабина за $\bigO(n^3)$.
\end{proposition}

\checkyourself{Почему нельзя взять $e = 2$? Почему нельзя взять $p = q$ (найдите $\varphi(p^2)$
и объясните, как противник разложит $N = p^2$)?}

\section{Стойкость RSA и факторизация}

Противник знает $(N, e)$ и $c$. Если он умеет раскладывать $N$, то вычисляет $\varphi(N)$ и
$d$ так же, как Алиса: \emph{факторизация ломает RSA}. Обратное "--- открытый вопрос:
\emph{задача RSA} (извлечь корень $e$-й степени по модулю $N$, не зная разложения) не
сложнее факторизации, но эквивалентность не доказана. Зато нахождение \emph{закрытого
ключа} эквивалентно факторизации.

\begin{proposition}\label{l5:phi-factor}
По $N = pq$ и $\varphi(N)$ числа $p$, $q$ находятся за полиномиальное время.
\end{proposition}

\begin{proof}
$\varphi(N) = N - (p + q) + 1$, поэтому $s = p + q = N - \varphi(N) + 1$ известно, и $p$, $q$ "---
корни уравнения $x^2 - s x + N = 0$:
\[
  p, q = \frac{s \mp \sqrt{s^2 - 4N}}{2}.
\]
Целочисленный квадратный корень вычисляется методом Ньютона за полиномиальное время.
\end{proof}

Для $N = 3233$, $\varphi(N) = 3120$: $s = 114$, $s^2 - 4N = 12996 - 12932 = 64$,
$p, q = (114 \mp 8)/2 = 53, 61$.

\taskexample{rsa-from-phi}

\begin{lemma}\label{l5:sqrt1}
Пусть $N = pq$, $p \ne q$ нечётные простые. Сравнение $x^2 \equiv 1 \pmod N$ имеет ровно
четыре решения: $\pm 1$ и ещё два, $\pm x_0$. Для любого \emph{нетривиального} корня
$x \not\equiv \pm 1$ число $\gcd(x - 1, N)$ "--- это $p$ или $q$.
\end{lemma}

\begin{proof}
По модулю простого $p$ из $p \mid (x - 1)(x + 1)$ следует $x \equiv \pm 1$ (лемма Евклида), и
$1 \not\equiv -1$, так как $p > 2$. По КТО решения по модулю $N$ "--- это четыре пары знаков
$(\pm 1 \bmod p, \pm 1 \bmod q)$. Если $x \not\equiv \pm 1 \pmod N$, то знаки различны,
например $x \equiv 1 \pmod p$, $x \equiv -1 \pmod q$; тогда $p \mid x - 1$, а $q \nmid x - 1$,
и $\gcd(x - 1, N) = p$.
\end{proof}

\begin{algorithm}[разложение $N$ по $e$ и $d$]\label{l5:factor-d}
\begin{algorithmic}[1]
\Require $N = pq$, $e$, $d$ с $ed \equiv 1 \pmod{\lambda(N)}$.
\Ensure $p$ или $q$.
\State $k \gets ed - 1$; записать $k = 2^t r$, $r$ нечётно
\While{делитель не найден}
  \State выбрать случайное $a \in [2, N - 2]$; \textbf{если} $g = \gcd(a, N) > 1$, \Return $g$
  \State $x \gets a^r \bmod N$
  \If{$x \ne 1$}
    \While{$x^2 \not\equiv 1 \pmod N$} $x \gets x^2 \bmod N$ \EndWhile
    \If{$x \ne N - 1$} \Return $\gcd(x - 1, N)$ \EndIf
  \EndIf
\EndWhile
\end{algorithmic}
\end{algorithm}

Это та же идея, что в тесте Миллера"--~Рабина (лекция~8 первого семестра): ищем
нетривиальный квадратный корень из единицы в последовательности $a^r, a^{2r}, \dots,
a^{2^t r} \equiv 1$.

\begin{theorem}\label{l5:d-factor}
Алгоритм разложения по $e$ и $d$ корректен, и для случайного $a$ одна итерация успешна с
вероятностью не меньше $1/2$. Ожидаемое число итераций не больше двух, т.\,е. по $(e, d)$
модуль $N$ раскладывается за ожидаемое полиномиальное время $\bigO(n^3)$.
\end{theorem}

\begin{proof}
Так как $\lambda(N) \mid k$, для $a \in \Znx{N}$ имеем $a^{2^t r} \equiv 1$, и внутренний
цикл останавливается. Возвращаемое $x$ удовлетворяет $x^2 \equiv 1$,
$x \not\equiv \pm 1$, и по лемме~\ref{l5:sqrt1} $\gcd(x - 1, N) \in \{p, q\}$.

Оценим вероятность неудачи. Обозначим через $2^{v_p(a)}$ порядок элемента $a^r$ в $\Znx{p}$
(это степень двойки, так как $(a^r)^{2^t} = 1$). Тогда
$a^{2^i r} \equiv 1 \pmod p$ ровно при $i \ge v_p(a)$, а при $i = v_p(a) - 1$ это значение
равно $-1 \bmod p$ (единственный корень из $1$, отличный от $1$). Аналогично для $q$.
Неудача означает $x_0 = a^r \equiv 1$ или появление $-1 \pmod N$ перед единицей, т.\,е.
в обоих случаях $v_p(a) = v_q(a)$; если же $v_p(a) \ne v_q(a)$, например $v_p < v_q$,
то $x = a^{2^{v_q - 1} r}$ сравнимо с $1$ по модулю $p$ и с $-1$ по модулю $q$ "--- успех.

Пусть $p - 1 = 2^{s} m$, $m$ нечётно. Так как $p - 1 \mid k$, то $m \mid r$, и порядок $a^r$
равен $2$-части порядка $a$. Группа $\Znx{p}$ циклична с образующей $g$; для $a = g^j$
порядок $a$ имеет $2$-часть $2^{s}$ ровно при нечётном $j$. Значит,
$\Pr[v_p(a) = s] = 1/2$, а остальные значения $v_p$ в сумме тоже имеют вероятность $1/2$:
каждое значение $v_p$ принимается с вероятностью не больше $1/2$. По КТО вычеты $a \bmod p$ и $a \bmod q$ независимы и равномерны, поэтому
\[
  \Pr[v_p(a) = v_q(a)] = \sum_{j} \Pr[v_q(a) = j]\,\Pr[v_p(a) = j]
  \le \max_j \Pr[v_p(a) = j] \le \frac12.
\]
Каждая итерация "--- одно возведение в степень и не более $t \le \log_2 k$ возведений
в квадрат, т.\,е. $\bigO(n)$ умножений по модулю $N$.
\end{proof}

Для $N = 3233$, $e = 17$, $d = 2753$: $k = 46800 = 2^4 \cdot 2925$. При $a = 2$ цепочка
$2^{2925} \equiv 1514$, $1514^2 \equiv 3232 \equiv -1$, $1$ "--- неудача; при $a = 3$:
$3^{2925} \equiv 2256$, $2256^2 \equiv 794$, $794^2 \equiv 1$, и $\gcd(793, 3233) = 61$.

\taskexample{rsa-from-d}

\begin{corollary}
Один модуль $N$ нельзя использовать для нескольких пользователей с разными парами $(e_i, d_i)$:
каждый из них, зная свою пару, раскладывает $N$ и находит чужие закрытые ключи.
\end{corollary}

\subsection{Атаки на «учебный» RSA}

\begin{proposition}[общий модуль]\label{l5:common-modulus}
Если одно и то же сообщение $m$ зашифровано на ключах $(N, e_1)$ и $(N, e_2)$ с
$\gcd(e_1, e_2) = 1$, то по $c_1$, $c_2$ сообщение восстанавливается без разложения $N$:
$m \equiv c_1^{u} c_2^{v} \pmod N$, где $u e_1 + v e_2 = 1$.
\end{proposition}

\begin{proof}
$c_1^u c_2^v \equiv m^{u e_1 + v e_2} = m$; одно из чисел $u, v$ отрицательно, соответствующий
множитель вычисляется как степень обратного элемента (если $c_i$ необратим, то
$\gcd(c_i, N)$ "--- делитель $N$).
\end{proof}

\begin{example}
$N = 3233$, $e_1 = 17$, $e_2 = 7$, $c_1 = 2790$, $c_2 = 1317$. Из $17 = 2 \cdot 7 + 3$,
$7 = 2 \cdot 3 + 1$ получаем $1 = 5 \cdot 7 - 2 \cdot 17$, т.\,е. $u = -2$, $v = 5$.
$c_1^{-1} \equiv 1788$, $1788^2 \equiv 2740$, $1317^5 \equiv 2000$, и
$m \equiv 2740 \cdot 2000 \equiv 65 \pmod{3233}$.
\end{example}

\begin{remark}[малые показатели и необходимость дополнения]
\begin{enumerate}
  \item Если $e = 3$ и $m < N^{1/3}$, то $c = m^3$ без приведения по модулю, и $m = \sqrt[3]{c}$
    вычисляется над целыми числами.
  \item \emph{Атака Хостада}: одно сообщение отправлено с $e = 3$ трём адресатам с попарно
    взаимно простыми модулями $N_1, N_2, N_3$. По КТО противник находит
    $C \equiv m^3 \pmod{N_1 N_2 N_3}$, а так как $m^3 < N_1 N_2 N_3$, то $C = m^3$ и
    $m = \sqrt[3]{C}$.
  \item «Учебный» RSA детерминирован (одинаковые сообщения дают одинаковые шифротексты,
    противник может проверить догадку о $m$) и \emph{мультипликативен}:
    $(m_1 m_2)^e = m_1^e m_2^e$, так что по $c$ можно получить шифротекст для $2m$, не зная $m$.
\end{enumerate}
Поэтому на практике сообщение перед шифрованием \emph{дополняется} случайными данными по
определённой схеме (OAEP, стандарт PKCS\,\#1), что устраняет все перечисленные атаки.
Обзор атак на RSA "--- в~\cite{boneh-rsa-attacks}.
\end{remark}

\checkyourself{Алиса и Боб случайно сгенерировали модули $N_A = p q_A$ и $N_B = p q_B$ с общим
простым $p$. Как Ева разложит оба модуля?}

\section{Алгоритм \texorpdfstring{$p - 1$}{p - 1} Полларда}

\subsection{Идея}

Пусть $p$ "--- неизвестный простой делитель $N$. По малой теореме Ферма $a^{p-1} \equiv 1
\pmod p$, а значит, $a^{M} \equiv 1 \pmod p$ для \emph{любого} кратного $M$ числа $p - 1$.
Тогда $p \mid a^M - 1$, и
\[
  p \mid \gcd(a^M - 1, N).
\]
Если при этом $a^M \not\equiv 1$ по модулю другого простого делителя $q$, то этот $\gcd$ "---
нетривиальный делитель $N$. Число $p$ мы не знаем, но если $p - 1$ раскладывается только на
\emph{малые} простые, то в качестве $M$ годится произведение всех малых степеней простых
"--- например $M = n!$ или $M = \lcm(1, 2, \dots, B)$ (Дж.~Поллард, 1974~\cite{pollard-pm1}).

\begin{definition}
Число $n$ называется \emph{$B$-гладким}, если все его простые делители не превосходят $B$, и
\emph{$B$-степенно-гладким}, если все степени простых $\ell^{\alpha} \parallel n$ не
превосходят $B$.
\end{definition}

Например, $1200 = 2^4 \cdot 3 \cdot 5^2$ "--- $5$-гладкое и $25$-степенно-гладкое. Число $n$
$B$-степенно-гладко тогда и только тогда, когда $n \mid \lcm(1, \dots, B)$, и
\[
  \lcm(1, 2, \dots, B) = \prod_{\ell \le B \text{ простое}} \ell^{\floor{\log_\ell B}}.
\]
Заметим также, что любое $B$-степенно-гладкое $n$ делит $B!$.

\begin{theorem}\label{l5:pm1}
Пусть $p \mid N$, $p - 1$ "--- $B$-степенно-гладкое, $M = \lcm(1, \dots, B)$ (или $M = B!$), и
$\gcd(a, N) = 1$. Тогда $p \mid g = \gcd(a^M - 1, N)$. Если, кроме того, $N$ имеет простой
делитель $q$, для которого $\operatorname{ord}_q a \nmid M$, то $1 < g < N$.
\end{theorem}

\begin{proof}
$p - 1 \mid M$, поэтому $a^M = (a^{p-1})^{M/(p-1)} \equiv 1 \pmod p$, т.\,е. $p \mid g$. Если
$\operatorname{ord}_q a \nmid M$, то $a^M \not\equiv 1 \pmod q$, т.\,е. $q \nmid g$ и $g \ne N$.
\end{proof}

\begin{algorithm}[$p - 1$ Полларда]\label{l5:pm1-alg}
\begin{algorithmic}[1]
\Require составное $N$, граница $B$, основание $a$ (например, $a = 2$).
\Ensure нетривиальный делитель $N$ или «неудача».
\State \textbf{если} $g = \gcd(a, N) > 1$, \Return $g$
\State $x \gets a$
\For{$j = 2, 3, \dots, B$}
  \State $x \gets x^{j} \bmod N$ \Comment{$x = a^{j!}$}
  \State $g \gets \gcd(x - 1, N)$
  \If{$1 < g < N$} \Return $g$ \EndIf
  \If{$g = N$} \Return «неудача: сменить $a$ или вернуться назад» \EndIf
\EndFor
\State \Return «неудача: увеличить $B$»
\end{algorithmic}
\end{algorithm}

На практике $\gcd$ вычисляют не на каждом шаге, а раз в несколько шагов (или по степеням
простых $\ell^{\floor{\log_\ell B}}$, $\ell \le B$, если $M = \lcm(1, \dots, B)$), запоминая
промежуточные значения $x$.

\begin{example}
$N = 1238231$, $a = 2$. Вычисляем $x_j = 2^{j!} \bmod N$ и $\gcd(x_j - 1, N)$:
\[
\begin{array}{c|rrrrrrrrr}
  j & 2 & 3 & 4 & 5 & 6 & 7 & 8 & 9 & 10\\ \hline
  x_j & 4 & 64 & 680213 & 137976 & 722863 & 1062990 & 774506 & 341300 & 1029258\\
  g_j & 1 & 1 & 1 & 1 & 1 & 1 & 1 & 1 & 1201
\end{array}
\]
$N = 1201 \cdot 1031$. Здесь $1200 = 2^4 \cdot 3 \cdot 5^2$ делит $10!$, а
$1030 = 2 \cdot 5 \cdot 103$ "--- нет (простой множитель $103$ появится только при $j = 103$).
\end{example}

\taskexample{pm1-factorial}

\subsection{Неудачи и как с ними бороться}

\begin{itemize}
  \item $g = 1$: ни для одного простого $p \mid N$ порядок $a$ не делит $M$. Скорее всего,
    у всех $p - 1$ есть простой множитель больше $B$; нужно увеличить $B$ (или применить
    другой метод).
  \item $g = N$: порядки $a$ по модулю \emph{всех} простых делителей $N$ разделили $M$
    одновременно "--- граница $B$ слишком велика для этого $N$ (или $N$ специально
    подобран). Возвращаемся к сохранённому промежуточному значению $x$ и добавляем
    множители показателя по одному: по степеням простых, затем по одному простому $\ell$.
    Если и это не помогает (оба порядка дописываются одним и тем же множителем $\ell$),
    меняем основание $a$.
\end{itemize}

\taskexample{pm1-backtrack}

\subsection{Сложность}

\begin{proposition}\label{l5:pm1-cost}
При $M = \lcm(1, \dots, B)$ вычисление $a^M \bmod N$ требует $\bigO(B)$ умножений по модулю
$N$, т.\,е. $\bigO(B \log^2 N)$ битовых операций (школьное умножение); вариант с $M = B!$ "---
$\bigO(B \log B)$ умножений.
\end{proposition}

\begin{proof}
Бинарный алгоритм возводит в степень $M$ за не более $2\log_2 M$ умножений, а
\[
  \ln \lcm(1, \dots, B) = \sum_{\ell \le B} \floor{\log_\ell B} \ln \ell = \psi(B) = \bigO(B)
\]
по теореме Чебышёва ($\psi(B) \sim B$ по асимптотическому закону распределения простых чисел;
без доказательства, см.~\cite[гл.~1]{crandall-pomerance}). Для $M = B!$
$\log_2 B! = \Theta(B \log B)$ по формуле Стирлинга. Вычисления $\gcd$ (их
$\bigO(B)$) стоят $\bigO(\log^2 N)$ каждое.
\end{proof}

Таким образом, время работы определяется не размером $N$, а размером \emph{наибольшей
степени простого} в $p - 1$ для наименее «неудобного» делителя $p$: метод $p - 1$ "---
\emph{специальный} алгоритм, эффективный для чисел особого вида. Для случайного простого $p$
вероятность того, что $p - 1$ окажется $B$-гладким, приблизительно равна $u^{-u}$, где
$u = \ln p / \ln B$ (теорема Дикмана"--~де~Брёйна, без доказательства,
см.~\cite[гл.~1]{crandall-pomerance}). Например, для $512$-битного $p$ и $B = 2^{40}$
имеем $u \approx 12{,}8$ и $u^{-u} \approx 10^{-14}$: для случайных простых RSA метод
практически бесполезен.

\begin{remark}[вторая стадия]
Чаще встречается ситуация, когда $p - 1 = (\text{$B_1$-гладкое}) \cdot Q$ с одним простым
$B_1 < Q \le B_2$, $B_2 \gg B_1$. После первой стадии $b = a^M$ перебирают простые
$Q_1 < Q_2 < \dots$ из $(B_1, B_2]$ и вычисляют $b^{Q_{i}} = b^{Q_{i-1}} \cdot b^{Q_i - Q_{i-1}}$,
где степени $b^{\delta}$ для малых чётных разностей $\delta$ вычислены заранее, накапливая
произведение $\prod (b^{Q_i} - 1) \bmod N$ и вычисляя $\gcd$ с $N$ лишь изредка: одно-два
умножения на каждое простое, так что $B_2$ может быть в $100$ и более раз больше $B_1$
\cite[разд.~5.4]{crandall-pomerance}.
\end{remark}

\subsection{Следствия для RSA и что дальше}

Если хотя бы для одного из простых $p$, $q$ число $p - 1$ гладкое, модуль $N$ раскладывается
методом Полларда. Поэтому стандарты RSA одно время требовали «сильных» простых: $p - 1$
должно иметь большой простой делитель $r$ (а $r - 1$ "--- тоже), и $p + 1$ "--- тоже (из-за
аналогичного метода $p + 1$ Уильямса, использующего группу порядка $p + 1$). Простейший
способ "--- \emph{безопасные} простые $p = 2p' + 1$ с простым $p'$. Для случайных простых
длиной $1024$ бит и более вероятность гладкости $p - 1$ ничтожна, и современные стандарты
ограничиваются выбором случайных простых достаточной длины.

Слабое место метода "--- фиксированная группа $\Znx{p}$ порядка $p - 1$: если $p - 1$ не
гладко, никакой выбор $a$ не поможет. Х.~Ленстра заменил её группой точек эллиптической
кривой $E(\FF_p)$, порядок которой меняется вместе с кривой в интервале
$[p + 1 - 2\sqrt p, p + 1 + 2\sqrt p]$; перебирая кривые, рано или поздно получаем гладкий
порядок. Это метод эллиптических кривых (лекция~10). Общие субэкспоненциальные методы "---
квадратичное решето "--- мы изучим в лекции~6.

\checkyourself{Почему в алгоритме $p - 1$ берут $a = 2$, а не $a = N - 1$? Что вернёт
алгоритм при $a = N - 1$?}

\checkyourself{Пусть $N = pq$, где $p - 1 = 2 \cdot 3 \cdot 5 \cdot 7$, а $q - 1 = 2 \cdot 11$.
Не позже какого $j$ вариант с $a^{j!}$ найдёт делитель $p$? Что, скорее всего, получится,
если сразу взять $M = \lcm(1, \dots, 11)$, и как тогда всё же найти делитель?}

\checkyourselfsection

\begin{furtherreading}
  \readitem[разд.~3.1--3.3]{hps}{теорема Эйлера, корни по модулю $pq$, RSA, вопросы стойкости}
  \readitem[разд.~3.5]{hps}{алгоритм $p - 1$ Полларда, гладкие числа}
  \readitem[разд.~5.4, 8.1.2]{crandall-pomerance}{метод $p - 1$ (с второй стадией); RSA}
  \readitem[гл.~1, 2; гл.~3, \S\,5]{vinogradov}{основная теорема арифметики, функция Эйлера,
    теоремы Эйлера и Ферма}
  \readitem[разд.~3.2.3, 8.2]{hac}{метод $p - 1$; RSA, его реализация и атаки}
  \readitem[гл.~7]{paar-pelzl}{RSA: генерация ключей, быстрое расшифрование по КТО, атаки}
  \readitem{boneh-rsa-attacks}{обзор атак на RSA}
  \readitem{msu-ent}{элементарная теория чисел: делимость, сравнения, функция Эйлера}
  \readitem[неделя~6]{coursera-boneh}{RSA как перестановка с секретом, атаки и дополнения}
  \readitem{rsa}{оригинальная статья RSA}
\end{furtherreading}
\end{lecture}
